\documentclass[11pt,a4paper]{article}

% ===== パッケージ =====
\PassOptionsToPackage{numbers,sort&compress}{natbib}

\usepackage[T1]{fontenc}
\usepackage{newtxtext,newtxmath}
\usepackage{amsmath,amsfonts,bm}
\usepackage{graphicx}
\usepackage{booktabs,multirow,array}
\usepackage{caption,subcaption}
\usepackage{algorithm2e}
\usepackage{relsize}
\usepackage[margin=2.5cm]{geometry}
\usepackage{setspace}
\usepackage[dvipsnames]{xcolor}
\usepackage[hyphens]{url}
\usepackage{natbib}  
\usepackage{multirow}

% ===== External References =====
\usepackage{xr}
\externaldocument{main}

% ===== Proposition Environment =====
\newtheorem{proposition}{Proposition}
\newtheorem{lemma}{Lemma}
\newtheorem{proof}{Proof}

\DeclareMathOperator*{\argmin}{arg\,min}

\renewcommand{\thefigure}{S\arabic{figure}}
\renewcommand{\thetable}{S\arabic{table}}
\renewcommand{\theequation}{S\arabic{equation}}

\renewcommand{\thesection}{S\arabic{section}}
\renewcommand{\topfraction}{0.8}     % ページ上部に float を詰める割合を緩める
\renewcommand{\bottomfraction}{0.8}  % ページ下部も少し緩める
\renewcommand{\textfraction}{0.2}   % テキストの最低割合を下げる

\begin{document}

% ===== Title and Authors =====
\title{Facet-induced connotation space: A library with rabbit holes where visitors seek serendipity}
\author{
    Tatsuya Kanatsu\textsuperscript{1}, Tohru Iwasaki\textsuperscript{2}, Hideaki Ishibashi\textsuperscript{1}, \\
    Kaori Yoshida\textsuperscript{1}, and Tetsuo Furukawa\textsuperscript{1} \\[3pt]
    \small \textsuperscript{1}Graduate School of Life Science and Systems Engineering, Kyushu Institute of Technology, Japan \\
    \small \textsuperscript{2}Data / AI System Division, ZOZO, Inc., Japan \\
    \small Corresponding author: furukawa@brain.kyutech.ac.jp
}
\date{}
\maketitle

\section{Use of Large Language Models}\label{sup:use_of_llms}
In preparing this manuscript, we used Large Language Models (LLMs) in the
following capacities, in accordance with the journal's guidelines.

\paragraph{Programming assistance.}
ChatGPT (OpenAI), Claude (Anthropic), and Cursor were used to assist in
developing the code for the computational experiments. Concretely, they
supported code generation and debugging for the data preprocessing, the
embedding extraction, the OT-FB learning, and the analysis and plotting
scripts.

\paragraph{Manuscript preparation.}
Claude (Anthropic) was used for translation assistance (converting the
authors' ideas from Japanese to English), language polishing (grammar,
sentence structure, and clarity), and consistency checking (terminology
throughout the manuscript).

\paragraph{Human oversight.}
All scientific content---conceptual frameworks, mathematical formulations,
experimental designs, results, and interpretations---originated entirely from
the authors. All LLM outputs were critically reviewed, and the authors retain
full responsibility for all claims and content in this manuscript.


\section{Acknowledgments}
We are deeply grateful to the reviewers. Their thorough and constructive
comments led us to rethink this work from its foundations, and the paper that
resulted is a far better one than the paper they received. We thank
Ms.\ Hirowatari for her kind assistance.

This research was supported in part by JSPS KAKENHI Grant Numbers 21K12061 and
26K14941. We also gratefully acknowledge the support of ZOZO, Inc.

\end{document}




